\documentclass[review]{elsarticle}

\usepackage{lineno,hyperref, graphicx, amsmath, esvect, amssymb, ragged2e, afterpage, placeins}
\usepackage[labelformat=simple]{subcaption}
\numberwithin{equation}{section}
\numberwithin{figure}{section}
\numberwithin{table}{section}
\usepackage{cleveref}
\renewcommand\thesubfigure{(\alph{subfigure})}
\modulolinenumbers[5]
      
\journal{Int J Appl Earth Obs Geoinformation}

\setcounter{secnumdepth}{4}

%\titleformat{\paragraph}
%{\normalfont\normalsize\it}{\theparagraph}{1em}{}
%\titlespacing*{\paragraph}
%{0pt}{3.25ex plus 1ex minus .2ex}{1.5ex plus .2ex}

%%%%%%%%%%%%%%%%%%%%%%%
%% Elsevier bibliography styles
%%%%%%%%%%%%%%%%%%%%%%%
%% To change the style, put a % in front of the second line of the current style and
%% remove the % from the second line of the style you would like to use.
%%%%%%%%%%%%%%%%%%%%%%%

%% Numbered
%\bibliographystyle{model1-num-names}

%% Numbered without titles
%\bibliographystyle{model1a-num-names}

%% Harvard
%\bibliographystyle{model2-names.bst}\biboptions{authoryear}

%% Vancouver numbered
%\usepackage{numcompress}\bibliographystyle{model3-num-names}

%% Vancouver name/year
%\usepackage{numcompress}\bibliographystyle{model4-names}\biboptions{authoryear}

%% APA style
\bibliographystyle{model5-names}\biboptions{authoryear}

%% AMA style
%\usepackage{numcompress}\bibliographystyle{model6-num-names}

%% `Elsevier LaTeX' style
%\bibliographystyle{elsarticle-num}
%%%%%%%%%%%%%%%%%%%%%%%

\tolerance=1
\emergencystretch=\maxdimen
\hyphenpenalty=10000
\hbadness=10000
\sloppy

\usepackage{scalerel}
\usepackage{tikz}
\usetikzlibrary{svg.path}

\definecolor{orcidlogocol}{HTML}{A6CE39}
\tikzset{
  orcidlogo/.pic={
    \fill[orcidlogocol] svg{M256,128c0,70.7-57.3,128-128,128C57.3,256,0,198.7,0,128C0,57.3,57.3,0,128,0C198.7,0,256,57.3,256,128z};
    \fill[white] svg{M86.3,186.2H70.9V79.1h15.4v48.4V186.2z}
                 svg{M108.9,79.1h41.6c39.6,0,57,28.3,57,53.6c0,27.5-21.5,53.6-56.8,53.6h-41.8V79.1z M124.3,172.4h24.5c34.9,0,42.9-26.5,42.9-39.7c0-21.5-13.7-39.7-43.7-39.7h-23.7V172.4z}
                 svg{M88.7,56.8c0,5.5-4.5,10.1-10.1,10.1c-5.6,0-10.1-4.6-10.1-10.1c0-5.6,4.5-10.1,10.1-10.1C84.2,46.7,88.7,51.3,88.7,56.8z};
  }
}

\newcommand\orcidicon[1]{\href{https://orcid.org/#1}{\mbox{\scalerel*{
\begin{tikzpicture}[yscale=-1,transform shape]
\pic{orcidlogo};
\end{tikzpicture}
}{|}}}}

\DeclareMathOperator{\sinc}{sinc}
\begin{document}

\begin{frontmatter}

\title{Pol-InSAR and PolSAR based Inversion Modelling for Snow Depth and SWE Estimation in the Northwestern Himalayan Watershed}
%\tnotetext[mytitlenote]{Fully documented templates are available in the elsarticle package on \href{http://www.ctan.org/tex-archive/macros/latex/contrib/elsarticle}{CTAN}.}

%% Group authors per affiliation:
%\author{Sayantan Majumdar\fnref{myfootnote}, Praveen K. Thakur, Ling Chang, Shashi Kumar}
%\address{Faculty ITC, University of Twente}
%\fntext[myfootnote]{Since 1880.}

%% or include affiliations in footnotes:
\author[itc,iirs]{Sayantan Majumdar \corref{corrauth}}
\cortext[corrauth]{Corresponding author}
\ead{monti.majumdar@gmail.com}
\author[iirs]{Praveen K. Thakur}
%\ead{praveen@iirs.gov.in}
\author[itc]{Ling Chang}
%\ead{ling.chang@utwente.nl}
\author[iirs]{Shashi Kumar}
%\ead{shashi@iirs.gov.in}

\address[itc]{Faculty of Geo-information Science and Earth Observation (ITC), University of Twente}
\address[iirs]{Indian Institute of Remote Sensing (IIRS), ISRO}

\begin{abstract}
Snow depth (SD) and Snow Water Equivalent (SWE) constitute essential physical properties of snow and find extensive usage in the hydrological modelling domain. However, the prominent influence of the hydrometeorological conditions present in the area of interest inhibits accurate large-scale measurement of the SD and SWE--- an ongoing research problem in the cryosphere paradigm. In this context, spaceborne synthetic aperture radar (SAR) systems benefit from global coverage at sufficiently high spatial resolutions. Previously, the copolar phase difference (CPD) method based on the X-band polarimetric SAR (PolSAR) technique has displayed promising results regarding the fresh snow depth (FSD) estimation. Still, this FSD inversion model has not been tested in the presence of extreme topographically varying conditions, such as the northwestern Himalayan belt. It is also susceptible to high volume scattering at X-band occurring from the increased snow grain sizes as a result of the standing (or old) snow formation driven by the temperature induced snow metamorphosis process. Hence, to model this volume decorrelation, the polarimetric SAR interferometry (Pol-InSAR) technique can be applied. In this work, the FSD and standing snow depth (SSD) are computed using the PolSAR CPD method and the single-baseline Pol-InSAR based hybrid Digital Elevation Model (DEM) differencing and coherence amplitude inversion model. To achieve this, the TerraSAR-X, TanDEM-X Coregistered Single look Slant range Complex (CoSSC) bistatic acquisition over Dhundi (situated in the Beas watershed, northwestern Himalayas, India) on January 8, 2016, is used. Due to the associated problems of model parameter tuning, complex topographical conditions, and limited ground-truth measurements, appropriate sensitivity analyses have been carried out for the parameter optimisation. Furthermore, the uncertainty sources are identified by performing a summer (June 8, 2017) and wintertime (January 8, 2016) comparative analysis of the study area which quantitatively highlights the changes in the percentages of the surface and volume scatterings. Evidently, the improved models display sufficiently high FSD and SSD accuracies of 94.83\% and 99.53\% respectively with the corresponding fresh SWE (FSWE) and standing SWE (SSWE) accuracies of 94.84\% and 99.48\% (measured using a 3$\times$3 window at Dhundi). Therefore, in summary, the overall outcome of this research showcases the practicability of these PolSAR and Pol-InSAR models in the context of the SD estimation over rugged terrains.
\end{abstract}

\begin{keyword}
Synthetic Aperture Radar, Copolar Phase Difference, Pol-InSAR, Snow Physical Properties, Sensitivity Analysis
\end{keyword}

\end{frontmatter}

\linenumbers

\section{Introduction}
\label{sec:Intro}

Snow depth (SD) and snow water equivalent (SWE) are two of the most important physical properties of snow and are extensively used in hydrological models that relate to snowmelt runoff and snow avalanche predictions \citep{Thakur2017}. While snow depth or snow height refers to the distance of the ground to the snow surface, SWE quantifies the amount of water present in a snowpack (layered snow formed by accumulation over time). Theoretically, SWE is defined as the product of snow depth and snow density and can be conceptualised as the amount of liquid water obtained owing to the instantaneous melting of
an entire snowpack \citep{Tedesco2015}. The accurate estimation of these two parameters is quite challenging depending upon the data availability and variety, mathematical model selection, and the hydrometeorological conditions of the area of interest. Hence, it is considered to be an important research element in the cryosphere paradigm \citep{Leinss2014, Leinss2015, Leinss2016, Conde2018}.

Due to the difficulties posed by in-situ or ground based measurements of snow depth and SWE in rugged terrains, remote sensing techniques coupled with adequately sampled (both in space and time domains) ground measurements are widely used to improve the quality of these estimated parameters over considerably large areas \citep{Takala2011}. Currently, LiDAR (Light Detection and Ranging) and spaceborne SAR are the most popular techniques used in the studies related to snow, ice and the cryosphere in general \citep{Deems2013, Leinss2014, Tedesco2015}. However, LiDAR can only be used to determine the height of the snow and cannot be used for measuring other physical properties such as snow density and snow wetness. In addition, the operating cost of LiDAR is sufficiently high and is also weather dependent \citep{Deems2013}. As a result, spaceborne SAR systems benefit from substantial coverage (globally available), cloud insensitivity, night-time operability and are extensively used to measure the snow physical properties sufficiently at high spatial resolutions \citep{Moreira2013, Thakur2012}.

The applicability of SAR systems for snow cover monitoring was discussed as early as 1977 \citep{Ulaby1977} wherein the snow backscatter coefficient was measured and was thereafter modelled for various frequencies, layers, and polarisations \citep{Zuniga1979}. It was shown that only very high microwave frequencies (Ku-band or higher) exhibit a significant dependence on SD or the SWE of dry or standing (deposited) snow \citep{Yueh2009}. However, lower frequencies (X-
band or below) penetrate through dry snow whereby the underneath frozen soil or ground primarily contributes to the radar backscatter signal. On the other hand, in case of moist snow (the transitional stage between dry and wet snow) and wet snow, the predominant scattering occurs from the snow volume and snow surface respectively due to the presence of water. Essentially, water, with its high dielectric constant, heavily modifies the dielectric properties of snow and effectively reduces the snow penetration capacity of the radar pulses \citep{Abe1990}. The radar backscattering mechanism for a typical snow covered area can be conceptualised from Figure \ref{fig:concept}. In principle, PolSAR and InSAR systems utilise these received target echoes for supporting various microwave remote sensing applications in the cryosphere domain.

\begin{figure}[htb]
    \centering
    \includegraphics[width=\textwidth]{Figures/Conceptual.png}
    \caption{Conceptual diagram displaying the radar backscattering mechanism in hilly terrains. Adapted from \cite{Thakur2012}.}
    \label{fig:concept}
\end{figure}

A polarimetric SAR system utilises the polarised radar echoes to obtain information about the specific scattering mechanism for a particular target. In essence, by using a coherent analysis which incorporates the phase of different polarimetric channels, it is possible to differentiate various scattering mechanisms \citep{Lee2009}. Nowadays, PolSAR based algorithms that work on the polarimetric backscatter signal have been widely adopted for various snow related applications such as the classification of dry and wet snow, measuring snow wetness and snow density \citep{Singh2017, Snehmani2010, Thakur2012, Thakur2017, Usami2016}. In this context, the roll-invariant entropy/anisotropy/alpha ($H/A/{\alpha}$) decomposition and Wishart classification have been successfully tested to classify different snow types as well as demarcate snow covered areas \citep{Cloude2010,Lee2009,Singh2014}. A few years back, the use of spaceborne PolSAR for snow height determination had been introduced, wherein the relationship between the copolar phase difference (CPD) and fresh snow depth (FSD) is quantitatively analysed by deriving a theoretical model \citep{Leinss2014}.

Interferometric SAR techniques find significant usage in the cryosphere domain and have been used to construct highly accurate digital elevation models (DEMs), measure dry snow depth and SWE in several studies \citep{Conde2018, Guneriussen2001, Leinss2015, Li2017, Liu2017}. The principle of SAR interferometry builds upon measured phase differences between radar images of the same area acquired at different temporal instances (repeat-pass) or different viewing geometries but same epoch (single-pass) \citep{Hanssen2001}. Still, the inherent problems of spatial and temporal decorrelations and atmospheric inhomogeneities are the primary limiting factors in the studies involving InSAR and its variant D-InSAR (Differential InSAR) \citep{Pepe2017}. These factors are responsible for inaccurate and low-coherence measurements, thereby leading to a potential decrease in the accuracy of the final results. In cryosphere research, the loss of coherence in InSAR is heavily influenced by the snow humidity, melting, and refreezing and is also susceptible to the variations in both spatial and temporal baselines. Although data assimilation algorithms like 3DVAR (three dimensional variation) and EnKF (Ensemble Kalman Filter) have been applied to the produced outputs of the SD inversion models for minimizing the effect of temporal decorrelation, the applicability and feasibility of such algorithms remains untested on varying data sets and study areas \citep{Liu2017}.

The Pol-InSAR technique works on the coherent combination of both PolSAR and InSAR observations, thereby enabling the interferogram generation in arbitrary transmit and receive channels \citep{Papathanassiou2001, Cloude2005, Cloude2010}. It has been widely used for estimating tree height in forested regions and can be effectively applied to natural or artificial volume scatterers including snow and ice \citep{Hajnsek2009, Kugler2015, Kumar2017, Papathanassiou2001}. In essence, the identification of different scattering processes (PolSAR) and the vertical profile sensitivity (InSAR) are unique to this technique. Therefore, the applicability of Pol-InSAR based SD retrieval could prove its potential in case of the standing snow depth (SSD) \citep{Thakur2012, Thakur2017}.

The prime focus of this research is to estimate the FSD and SSD using PolSAR and Pol-InSAR respectively. In addition, the corresponding fresh SWE (FSWE) and standing SWE (SSWE) are to be determined, for which the respective snow densities need to be known. Essentially, the study involves sensitivity analysis of the various model parameters along with the identification of the different uncertainty sources present in the chosen geographical area. 

This manuscript is compartmentalised into five sections each consisting of several subsections. It starts with an introductory discussion in section \ref{sec:Intro}. Thereafter, the study area and the datasets including the required  software are specified in section \ref{sec:study}. From section \ref{sec:method} onwards the methodology and results are discussed. Finally, the relevant conclusions and recommendations are put forward in section \ref{sec:conc}.

\section{Study Area, Datasets, and Software}
\label{sec:study}

\subsection{Chosen Study Area}
\subsubsection{Geographical Situation}
\label{sssec:geo}
The Beas river watershed near Manali, India is part of the north-western Himalayas. Naturally, steep slopes and dense forests are prominent in this region. The elevation typically varies from nearly 2500 m to more than 5000 m in some places as observed in the reference ALOS PALSAR DEM (Figure \ref{fig:overview}). 

\begin{figure}[htb]
    \centering
    \includegraphics[width=\textwidth]{Figures/Overview.png}
    \caption{Overview map of the study area showing the ALOS PALSAR DEM. The original DEM of 12.5 m spatial resolution (generated in 2011) has been resampled to 3 m using bilinear interpolation \citep{Wu2008} to match the high resolution SAR data. Moreover, the vertical resolution as per the product specification is 5 m.}
    \label{fig:overview}
\end{figure}

In this work, a small region ($\sim$ 96 km\textsuperscript{2}) of the Beas basin is chosen which starts a few kilometres uphill from Dhundi up to Kothi (shown in Figure \ref{fig:overview}). These areas receive substantial seasonal snowfall which begins in December and lasts till late March. However, the cold, dry season usually commences from late September or early October. The coldest period is in January during which the temperatures reach a daily minimum of -15$^\circ$C on an average. The summers are mild to occasionally warm with June being the hottest month (mean and maximum temperatures of 20$^\circ$C and 30$^\circ$C respectively are common). Apart from this, significant rainfall occurs between late June and September (monsoon season) with August receiving the maximum precipitation \citep{Majumdar2019, Thakur2012}.

\subsubsection{Field Visit}
\label{sssec:field}
Intensive fieldwork had been conducted from October 14-21, 2018 in the Dhundi and Kothi areas where several Differential Global Positioning System (DGPS) measurements were acquired using the Leica Viva GS 10 \citep{LeicaGeosystemsAG2012} with adequate horizontal positional accuracies ($\sim$7 cm) \citep{Majumdar2019}. Due to the complex terrains, most of the DGPS readings had been obtained through the kinematic mode \citep{Luo2014}. However, in some of the convenient places such as the Dhundi base station and near the Kothi Automatic Weather Station (AWS), the static mode was used \citep{LeicaGeosystemsAG2012}. Eventually, elevation information from these DGPS points have been compared with the ALOS PALSAR DEM, the details of which are provided in section \ref{sec:res}. Furthermore, the manual snow readings from 2014-2018 (snow depth, density, weather profile and other relevant data) which are maintained by the security personnel daily at Dhundi had been pagewise photographed using a smartphone camera. In order to properly understand and visualise the characteristics of the study area, selected field photographs and their brief description are shown from figures \ref{subfig:gps}-\ref{subfig:stations}.

\begin{figure}[!t]
    \centering
    \begin{subfigure}[t]{0.49\textwidth}
        \raisebox{-\height}{\includegraphics[width=\textwidth]{Figures/Field/me_gps.jpg}}
        \caption{DGPS positional accuracy checking}
        \label{subfig:gps}
    \end{subfigure}
    \hfill
    \begin{subfigure}[t]{0.49\textwidth}
        \raisebox{-\height}{\includegraphics[width=\textwidth]{Figures/Field/base.png}}
        \caption{Leica DGPS base}
        \label{subfig:base}
    \end{subfigure}
    %%%%%%%%%%%%%%%%%%%%%%%%%%%%%%%%%%%%second row
    \begin{subfigure}[t]{0.49\textwidth}
        \raisebox{-\height}{\includegraphics[width=\textwidth]{Figures/Field/beas.jpg}}
        \caption{Beas river}
        \label{subfig:beas}
    \end{subfigure}
    \hfill
    \begin{subfigure}[t]{0.49\textwidth}
        \raisebox{-\height}{\includegraphics[width=\textwidth]{Figures/Field/landscape.jpg}}
        \caption{Landscape and human settlements}
        \label{subfig:landscape}
    \end{subfigure}
    %%%%%%%%%%%%%%%%%%%%%%%%%%%%%%%%%%%third row
    \begin{subfigure}[t]{0.49\textwidth}
        \raisebox{-\height}{\includegraphics[width=\textwidth]{Figures/Field/mountain.jpeg}}
        \caption{Mountains and forests}
        \label{subfig:mountain}
    \end{subfigure}
    \hfill
    \begin{subfigure}[t]{0.49\textwidth}
        \raisebox{-\height}{\includegraphics[width=\textwidth]{Figures/Field/stations.jpg}}
        \caption{Weather instruments}
        \label{subfig:stations}
    \end{subfigure}
    \caption{Dhundi field photographs showing the varying topographic features present in the surrounding area.}
    \label{fig:field}
\end{figure}

\subsection{Datasets Used}
\label{ssec:data}
Overall twelve Coregistered Single look Slant range Complex (CoSSC) TerraSAR-X (TSX)/TanDEM-X (TDX) bistatic X-band SAR images acquired between December 2015 and August 2017 in stripmap (SM) mode are available over this study area \citep{Balss2012}. In total, there are six Quad-pol data pairs, from which the descending orbital pass acquisition at 00:53 hrs Universal Time Coordinated (UTC), January 8, 2016, has been selected considering the occurrence of fresh snowfall before, during and after the satellite flyby. Moreover, the perpendicular baseline ($B_\bot$) and ambiguity height ($h_{2\pi}$) for this data are 96.34 m and 63.18 m respectively. 

Additionally, the in-situ snow physical parameters’ data (standing and fresh snow depths, snow density) along with the relevant weather data had been transferred to a PostgreSQL database (DB) \citep{PostgreSQL2019} from the photographs of the manual recordings through spreadsheets. Apart from this, the high frequency data (two-minute interval measurements) obtained from the snowpack analyser (SPA) device (installed at Dhundi) had been downloaded and were added to the database as a separate table. Accordingly, the SSDs at 06:22 hrs (00:52 hrs UTC) Indian Standard Time (IST) on January 7, 2016, and 06:22 hrs January 8, 2016 morning were 36.2 cm and 54.9 cm respectively signifying a heavy fresh snowfall event of 18.7 cm within 24 hrs. The manual recordings also showed an FSD of 18 cm on January 8, 2016 morning though the exact measurement time is unspecified in the record book. Apart from this, a forest mask used in the earlier studies of this area \citep{Thakur2012, Thakur2017} has been obtained from the Water
Resources Department (WRD), Indian Institute of Remote Sensing (IIRS).

\subsection{Software}

The Sentinel Application Platform (SNAP) 6.0.5 \citep{ESA2018} has been used for basic SAR processing. In addition, the FSD and SSD inversion models have been implemented using Python 3 wherein PyCharm Community Edition 2018.1 \citep{JetBrains2018} was used as the coding environment. Moreover, the final snow depth maps have been prepared using QGIS 2.18 \citep{QGIS2016}. Furthermore, some of the computationally intensive tasks have been carried out using the High-Performance Computing (HPC) infrastructure installed at IIRS.

\section{Methodology}
\label{sec:method}

This section deals with the methodological framework which has been followed to generate the SD and SWE results. In order to briefly put the overall workflow, a flowchart is shown in Figure \ref{fig:method_overview} which highlights the main process blocks. Here, the preprocessing steps are discussed in section \ref{ssec:pre}. Moreover, the PolSAR CPD and Pol-InSAR based approaches used for the FSD and SSD estimation respectively are individually addressed in sections \ref{ssec:fsd} and \ref{ssec:ssd}. Finally, the uncertainty assessment, validation and sensitivity tasks are described in section \ref{ssec:vus}.

\begin{figure}[htb]
    \centering
    \includegraphics[width=\textwidth]{Figures/Methods/Overview_Method.png}
    \caption{Overview of the main processing blocks.}
    \label{fig:method_overview}
\end{figure}

\subsection{Data Preprocessing}
\label{ssec:pre}

Since the SAR dataset is already coregistered, separate coregistration step has not been performed. In case of the FSD estimation model, the geocoded or terrain-corrected data (3 m spatial resolution) consists of the HH and VV scenes along with the local incidence angle (LIA) computed from the ALOS PALSAR DEM. As for the Pol-InSAR scenario, all the SAR channels, i.e., HH, HV, VH and VV along with the LIA are present in the geocoded data. It should be noted that, for the Pol-InSAR, processing both the master, TDX (master) and TSX (slave) images are required to generate the interferogram. However, the FSD estimation model can be used using any one of these images, though the average of the TDX and TSX CPDs can potentially improve the signal-to-noise ratio (SNR) \citep{Leinss2014}.

\subsection{CPD based Fresh Snow Depth Estimation}
\label{ssec:fsd}
\subsubsection{CPD Computation}

The FSD is estimated using the CPD method developed by \cite{Leinss2014}. At first, $\phi_{CPD}$ for the TDX data ($\phi_{CPD,TDX}$) acquired on January 8, 2016, is computed using Eq. \eqref{seq:cpd} and then an ensemble averaging operation is applied over a 21$\times$21 window \citep{Majumdar2019}. Similarly, $\phi_{CPD}$ for the TSX data ($\phi_{CPD,TSX}$) is calculated following which the average CPD, $\overline{\phi_{CPD}}$ is obtained using Eq. \eqref{seq:acpd}. Here, $\Im$ and $\Re$ denote the imaginary and real parts of the complex scattering matrices $S_{VV}$ and $S_{HH}$ respectively.

\begin{subequations}
    \centering
    \begin{equation}
        \label{seq:cpd}
        \phi_{CPD} = \phi_{VV} - \phi_{HH} = \left\langle\underbrace{\tan^{-1}\left(\frac{\Im\left(S_{VV}\right)}{\Re\left(S_{VV}\right)}\right)}_{\phi_{VV}} - \underbrace{\tan^{-1}\left(\frac{\Im\left(S_{HH}\right)}{\Re\left(S_{HH}\right)}\right)}_{\phi_{HH}}\right\rangle
    \end{equation}
    
    \begin{equation}
        \label{seq:acpd}
        \overline{\phi_{CPD}} = \frac{\phi_{CPD,TDX} + \phi_{CPD,TSX}}{2}
    \end{equation}
\end{subequations}

 $\phi_{CPD}$ can be alternatively defined as the phase angle of the complex copolar coherence, $\widetilde{\gamma_c}$ (since $\gamma$ is the standard notation for coherence amplitude, $\widetilde{\gamma}$ is used for the complex coherence), defined in Eq. \eqref{eq:2}. In this case, the copolar coherence amplitude ($\gamma_c = \lvert\widetilde{\gamma_c}\rvert$) is a measure of the radar backscattering mechanism where low values close to zero (ideally $\gamma_c =$ 0) indicate the presence of volume scattering and high values (ideally $\gamma_c =$ 1) represent surface scattering \citep{Lee2009, Leinss2014, Singh2014}. It should be noted that only the CPD computed from side-looking radar systems is able to describe a target having dielelectrically anisotropic microstructure as in the case of snow \citep{Leinss2014}. Also, in Eq. \eqref{eq:2}, $j$ is the imaginary unit. 

\begin{equation}
    \label{eq:2}
    \widetilde{\gamma_c} = \frac{\left\langle S_{VV}S_{HH}^*\right\rangle}{\sqrt{\left\langle S_{VV}S_{VV}^*\right\rangle\left\langle S_{HH}S_{HH}^*\right\rangle}} = \gamma_ce^{j\phi_{CPD}}, \gamma_c \in [0, 1]
\end{equation}
\subsubsection{Calculating the Depolarisation Factors}

In order to model this snow anisotropy, an ice particle needs to be associated with a specific shape. It has been observed that fresh snow and old snow exhibit horizontally aligned (oblate) and vertically aligned (prolate) spheroidal structures respectively \citep{Leinss2014}. Moreover, a shape parameter, known as the depolarisation factor, also has to be considered in this context \citep{Leinss2014, Sihvola1999}. In principle, a single spheroidal particle is characterised by three dipoles corresponding to the three orthogonal axes ($a_x$ , $a_y$ , and $a_z$) represented using a 3D ($x$, $y$, $z$) Cartesian coordinate system. This is depicted in Figure \ref{fig:prolate}, where the prolate shaped ice grain is linked with the radar reference frame ($h$, $k$, $v$) following the radar backscattering alignment (BSA) convention, $k$ being the propagation vector, $h$ and $v$ are the wave components of the horizontal and vertical polarisations respectively. Also, $\theta$ is the mean incidence angle with respect to the surface normal \citep{Leinss2014, Parrella2013}. 

\begin{figure}[htb]
    \centering
    \includegraphics[width=\textwidth]{Figures/Prolate_New.png}
    \caption{Orientation of a single prolate ice particle linked with the radar reference frame. Adapted from \cite{Leinss2014}.}
    \label{fig:prolate}
\end{figure}

So, by fixing a particle shape, the three depolarisation factors, $N_i$ ($\forall i \in \{x, y, z\}$), can be obtained by
solving the surface integral ($s$ is the ellipsoidal surface) as shown in Eq. \eqref{eq:3}.

\begin{equation}
    \label{eq:3}
	N_i =\frac{a_xa_ya_z}{2}\int_0^\infty\frac{\mathrm{d}s}{(s+a_i^2)\sqrt{(s+a_x^2)(s+a_y^2)(s+a_z^2)}}
\end{equation}
where, $N_x + N_y + N_z = 1$

For a perfectly spherical shape, all three depolarisation factors are equal to 1/3. The two other special cases include disk (1, 0, 0) and needle (0, 1/2, 1/2). In cases of prolate and oblate spheroids, the closed form expressions are already available as shown in Eq. \eqref{eq:4} \citep{Sihvola1999}. Here, the shape is dependent on the axial ratio ($a_x/a_z$) which is used for calculating the prolate eccentricity, $e_1 = \sqrt{1 - (a_x/a_z)^2}$ , and oblate eccentricity, $e_2 = \sqrt{(a_x/a_z)^2 - 1}$ respectively, i.e., for prolate, $a_x/a_z < 1$, whereas for oblate, it
is the reverse. However, for general ellipsoids having different axes, the above surface integration needs to be explicitly solved.

\begin{equation}
    \label{eq:4}
	N_z = \begin{cases} 
          \frac{1 - e_1^2}{2e_1^3}\left(\ln{\frac{1 + e_1}{1 - e_1} - 2e_1}\right) &, a_z > a_x = a_y \\
          \frac{1 + e_2^2}{e_2^3}\left(e_2 - \tan^{-1}e_2\right) &, a_z < a_x = a_y
       \end{cases}
\end{equation}

\subsubsection{Computing the Snow Refractive Indices}

Evidently, the Maxwell-Garnett theory related to electromagnetic mixing models can be applied to a medium (here snow) consisting of both air and ice which are having relative permittivities (real part), $\epsilon_{air}$ and $\epsilon_{ice}$ respectively. Therefore, the effective permittivity of this mixed medium, $\epsilon_{eff, i}$, is anisotropic and given by Eq. \eqref{eq:5} \citep{Leinss2014, Sihvola1999}. Here, the particle volume fraction ($f_{vol}$) is dependent on the fresh snow density ($\rho_f$)and ice density ($\rho_{ice}$). 

\begin{equation}
    \label{eq:5}
	\epsilon_{eff, i} = \epsilon_{air}\left[1 + f_{vol}\frac{\epsilon_{ice} - \epsilon_{air}}{\epsilon_{air} + \left(1 - f_{vol}\right)N_i\left(\epsilon_{ice} - \epsilon_{air}\right)}\right]
\end{equation}
where, $f_{vol} = \rho_f / {\rho_{ice}}, \rho_{ice} = 0.917 \text{ g/cm\textsuperscript{3}}, \epsilon_{air} = 1.00059, \epsilon_{ice} = 3.179, \text{ and } i \in \{x, y, z\}$

Furthermore, the refractive indices of this birefringent (or birefractive) medium, $n_H$ and $n_V$ corresponding to the HH and VV polarisations respectively, are dependent on this anisotropic effective permittivity \citep{Leinss2014}. In addition, since the snow anisotropy is assumed to be oriented along the Earth’s gravitational field, $n_H$ remains constant whereas $n_V$ is dependent on the incidence angle $\theta$ as given by Eq. \eqref{seq:nh} and Eq. \eqref{seq:nv} respectively \citep{Leinss2016}. Also, the imaginary part of the effective permittivity is negligible in the case of dry snow (fresh snow is also dry), and therefore, it is not used in the model developed by \cite{Leinss2014}.

\begin{subequations}
    \centering
    \begin{align}
        \label{seq:nh}
        n_H^2 &= \epsilon_{eff, x}\\
        \label{seq:nv}
        n_V^2 &= \epsilon_{eff, y}\cos^2\theta + \epsilon_{eff, z}\sin^2\theta 
    \end{align}
\end{subequations}
where, $\epsilon_{eff, x}$ , $\epsilon_{eff, y}$ , and $\epsilon_{eff, z}$ represent the effective permittivities of fresh snow in $x$, $y$, and $z$ directions of a 3D Cartesian co-ordinate system \citep{Leinss2014}.

\subsubsection{FSD and FSWE Computation}

Once all these aforementioned parameters are calculated, the CPD based inversion model given by Eq. \eqref{eq:7} is applied to estimate the depth of fresh snow, denoted by $\Delta{Z_f}$ \citep{Majumdar2019, Leinss2014, Leinss2016}. In this equation, -1 is introduced as per the BSA convention which is followed for all radar systems. Here, $\lambda_0$ is the radar wavelength and $\Delta\zeta$ is the relative path length difference which is dependent on $\epsilon_{eff, i}$, $\theta$, and $\rho_f$ \citep{Leinss2016}. Moreover, the horizontally aligned microstructure of fresh snow reduces the propagation speed for the HH channel and hence, in this case, $n_H > n_V$ always holds. However, for recrystallised snow having vertically aligned structures, the reverse condition is true \citep{Leinss2016}. Also, the LIA ($\theta_l$) is used instead of the mean incidence angle ($\theta$) to consider the effect of the terrain slope.

\begin{equation}
    \centering
    \label{eq:7}
    \Delta{Z_f} = \left\langle(-1)\frac{\lambda_0\phi_{CPD}}{4\pi\Delta{\zeta}}\right\rangle
\end{equation}
\begin{center}where, $\Delta{\zeta} = \sqrt{n^2_V - \sin^{2}\theta_l} - \sqrt{n^2_H - \sin^{2}\theta_l}, \phi_{CPD} > 0, \text{ and } n_H > n_V$\end{center}

Here, the depolarisation factors are calculated by setting the axial ratio, $a_x/a_z = 1.5$ in
Eq. \eqref{eq:4} and choosing the snow particle shape as an oblate \citep{Majumdar2019}. After this, the anisotropic effective permittivities are computed using Eq. \eqref{eq:5}. Finally, the FSD is obtained from Eq. \eqref{eq:7} wherein an ensemble averaging filter of size 65$\times$65 is applied. The FSWE is obtained by multiplying the FSD with the fresh snow density (i.e, FSWE $ = \rho_f\Delta{Z_f}$)Also, the fresh snow density ($\rho_f = $ 0.07 g/cm\textsuperscript{3}) which is manually measured at Dhundi, is kept constant for the entire study area along with the copolar coherence threshold, $\tau_c = $ 0 ($\tau_c \in$ [0, 1]), i.e., no thresholding has been applied on $\gamma_c$ computed using Eq. \eqref{eq:2}, but the provision for it is built-in to the implementation. Moreover, as per the TSX/TDX metadata, the radar wavelength, $\lambda_0 \approx$ 3.11 cm. In this context, the adopted workflow is depicted in Figure \ref{fig:fsd_method}.

\begin{figure}[htb]
    \centering
    \includegraphics[width=\textwidth]{Figures/Methods/FSD_Method.png}
    \caption{FSD and FSWE estimation workflow using PolSAR CPD.}
    \label{fig:fsd_method}
\end{figure}

\subsection{Pol-InSAR based Standing Snow Depth Estimation}
\label{ssec:ssd}

Standing or old snow refers to the deposited snow on the ground which has accumulated over time \citep{Reynolds1983}. Typically, old snow due to the presence of impurity and temperature-gradient induced recrystallisation process consists of snow particles larger than the microwave wavelength and results in volume scattering \citep{Leinss2016, Riche2013}. This volume decorrelation can be quantitatively analysed with the help of the Pol-InSAR technique \citep{Cloude2010} to obtain the volumetric SSD ($\Delta{Z_s}$). 

\subsubsection{Single-baseline Pol-InSAR Specifics}

The single baseline Pol-InSAR algorithm works on the basis of the complex coherence, $\widetilde{\gamma}\left(\vv{w_1}, \vv{w_2}\right)$, defined in Eq. \eqref{seq:coh} where $I_i\left(\vv{w_1}, \vv{w_2}\right)$ denotes the $i$\textsuperscript{th} pixel coordinate value of the wrapped Pol-InSAR interferogram, $I\left(\vv{w_1}, \vv{w_2}\right)$ obtained in Eq. \eqref{seq:ifg}. This interferogram is calculated from Eq. \eqref{seq:s1} and Eq. \eqref{seq:s2} where the coregistered master ($s_1$) and slave ($s_2$) images are acquired at a given polarisation vector, ($\vv{w}$) respectively. Here, the weight vectors, $\vv{w_1}$ and $\vv{w_2}$ are selected by the user based on the scattering mechanisms at ends 1 and 2 of the interferometric baseline. If $\vv{w_1} = \vv{w_2}$, $\widetilde{\gamma}\left(\vv{w_1}, \vv{w_2}\right)$ can be alternatively specified as $\widetilde{\gamma}\left(\vv{w_1}\right)$ \citep{Cloude2005, Cloude2010}. Moreover, $L$ is the total number of pixels averaged in the range and azimuth directions which can be replaced by the ensemble averaging operation following the statistical ergodicity assumption \citep{Hanssen2001, Hoen2000, Kugler2015, Kumar2017, Papathanassiou2001}. Additionally, $\phi_{flat}^w \in [0, 2\pi)$ is the wrapped flat-earth phase obtained from the estimated absolute flat-earth phase, $\phi_{flat}$ and has to be removed from $I\left(\vv{w_1}, \vv{w_2}\right)$ as shown in Eq. \eqref{seq:ifg}. Also, the calculation of the generalised weight vector, $\vv{w}$ is given by Eq. \eqref{seq:weight}.

\begin{subequations}
    \centering
    \begin{align}
        \label{seq:coh} 
        & \widetilde{\gamma}\left(\vv{w_1}, \vv{w_2}\right) = \frac{\sum_{i=1}^LI_i\left(\vv{w_1}, \vv{w_2}\right)}{\sqrt{\sum_{i=1}^L\left\lvert s_{1i}\left(\vv{w_1}\right)\right\rvert^2\sum_{i=1}^L\left\lvert s_{2i}\left(\vv{w_2}\right)\right\rvert^2}}, \left\lvert\widetilde{\gamma}\left(\vv{w_1}, \vv{w_2}\right)\right\rvert  \in [0, 1] \\
        \label{seq:ifg} 
        & I\left(\vv{w_1}, \vv{w_2}\right) = s_1\left(\vv{w_2}\right)s_2^*\left(\vv{w_2}\right)e^{-j\phi_{flat}^w} \\
        \label{seq:s1} 
        & s_1 = w_1^1\frac{s_{hh}^1 + s_{vv}^1}{\sqrt{2}} + w_1^2\frac{s_{hh}^1 - s_{vv}^1}{\sqrt{2}} + w_1^3\sqrt{2}s_{hv}^1 \\
        \label{seq:s2} 
        & s_2 = w_2^1\frac{s_{hh}^2 + s_{vv}^2}{\sqrt{2}} + w_2^2\frac{s_{hh}^2 - s_{vv}^2}{\sqrt{2}} + w_2^3\sqrt{2}s_{hv}^2 \\
        \label{seq:weight}
        & \vv{w} = \begin{bmatrix}
                    w^1 & w^2 & w^3
                    \end{bmatrix}^T = \begin{bmatrix}
                    \cos\alpha & \sin\alpha \cos{\beta{e}^{j\delta}} & \sin\alpha\sin{\beta{e}^{j\mu}}
                    \end{bmatrix}^T
    \end{align}
\end{subequations}
where, $s_{pp}^1$ and $s_{pp}^2$ correspond to the master (1) and slave (2) images respectively, $pp \in \{hh, hv, vv\}$,
and $*$ denotes the complex conjugate operator.

In this case, the parameters, scattering alpha angle ($\alpha$), target orientation angle ($\beta$), phase terms ($\delta$ and $\mu$), are chosen according to the selected polarisation given by Table \ref{table:1}. Here, LL, LR and RR correspond to
the left circular, left-right circular and right circular polarisations \citep{Cloude2010}. However, it is possible to
optimise these parameters specific to the data, the details of which are provided by \cite{Cloude2010}.

\begin{table}[ht]
\centering
\caption{Pol-InSAR scattering mechanisms \citep{Cloude2005}}
\label{table:1}
\begin{tabular}{c c c c c}
\hline
\textbf{Polarisation Selection} & \boldmath$\alpha(^\circ)$   & \boldmath$\beta(^\circ)$     & \boldmath$\delta(^\circ)$ & \boldmath$\mu(^\circ)$       \\ \hline
HH                     & 45 & 0   & 0          & 0   \\ 
HV                     & 90 & 90  & 0          & 0   \\ 
VV                     & 45 & 180 & 0          & 0   \\ 
HH+VV                  & 0  & 0   & 0          & 0   \\ 
HH-VV                  & 90 & 0   & 0          & 0   \\ 
LL                     & 90 & 45  & 0          & 90  \\ 
LR                     & 0  & 0   & 0           & 0   \\ 
RR                     & 90 & 45  & 0          & -90 \\ \hline
\end{tabular}
\end{table}

\subsubsection{Height Inversion Algorithm Details}
In this study, the modified (also improved) hybrid DEM differencing and coherence amplitude based Pol-InSAR volumetric height inversion model as given by Eq. \eqref{seq:hybrid_dem} is used for the SSD estimation. Firstly, the volume scattering dominant channels, HV and VH, are averaged to fully utilise the quad-pol data \citep{Cloude2005}. Next, the Pol-InSAR interferogram, $I\left(\vv{w_v}\right)$ has been computed using Eq. \eqref{seq:ifg} wherein, the weight vector, $\vv{w_v}$ is obtained from Table \ref{table:1} for the HV polarisation. Thereafter, the complex volume coherence, $\widetilde{\gamma}\left(\vv{w_v}\right)$, is calculated from Eq. \eqref{seq:coh} with $L$ = 3. Similarly, the complex surface or ground coherence, $\widetilde{\gamma}\left(\vv{w_s}\right)$, is computed by choosing $\vv{w_s}$ as the HH-VV weight vector (Table \ref{table:1}).  

Moreover, the actual vertical wavenumber, $k_z$ , when varied with the LIA, is in the order of 0.1 rad/m with the ambiguity height, $h_{2\pi} = 2\pi/k_z \approx$ 63.18 m, $\lambda_0 \approx$ 3.11 cm and $m$ = 1 (single-pass acquisition). Since the maximum height of the distributed volume scatterer (in this case, standing snow), $\Delta{Z_{s,max}}$, should be similar to $h_{2\pi}$ \citep{Kugler2015, Hajnsek2009, Kumar2017}, $k_z$ has to be rescaled to an optimum range for effectively estimating the SSD. Hence, the modified vertical wavenumber, $k_z^\prime$ , is given by Eq. \eqref{seq:vn} where $\eta^\prime$ is a free scaling parameter which has to be set according to the known $\Delta{Z_{s,max}}$ in the study area. Here, $h_{2\pi}^\prime$ is the scaled height of ambiguity which like that of $h_{2\pi}$ determines the height changes in modulo $2\pi$ \citep{Hanssen2001}. Also, $\mathbb{R}_{>0}^+$ denotes the set of all positive real numbers in the interval $(0, \infty)$. In this work, due to the limited ground-truth data availability and the subsequent ensemble
averaging operation (window size of 21$\times$21) on $k_z^\prime$ , $\Delta{Z_{s,max}} = $ 12 m has been assumed for which $\eta^\prime = $ 5 is used.

Apart from this, the function $\arg$ is defined in the interval $[0, 2\pi)$ and the parameter $m$ is set to 1 for bistatic acquisition and 2 in the monostatic case. Also in Eq. \eqref{seq:vn}, $\Delta\theta$ is the change in the incidence angle occuring due to the spatial baseline, $\theta_l$ is the LIA, and $\lambda_0$ is the radar wavelength as before in section \ref{ssec:fsd} \citep{Cloude2010, Kugler2015}.

\begin{subequations}
    \centering
    \begin{equation}
        \label{seq:hybrid_dem}
        \Delta{Z_s} = \frac{\arg\left(\widetilde{\gamma}\left(\vv{w_v}\right)e^{-j\phi_{topo}^w}\right)}{k_z^\prime} + \eta\frac{\sinc_\pi^{-1}\left(\gamma\left(\vv{w_v}\right)\right)}{k_z^\prime}, \eta\in[0, 1]
    \end{equation}
    \begin{flushleft}\text{where,}\end{flushleft}
    \begin{equation}
        \label{seq:vn}
        k_z^\prime = \left\langle\eta^\prime\frac{m\Delta\theta}{\lambda_0\sin\theta_l}\right\rangle, \eta^\prime\in\mathbb{R}_{>0}^+\mid\Delta{Z_{s, max}} \approx h_{2\pi}^\prime = 2\pi/k_z^\prime
    \end{equation}
\end{subequations}

Subsequently, the volume and surface coherences are then used to estimate the wrapped ground phase, $\phi_{topo}^w \in [0, 2\pi)$, from Eq. \eqref{eq:topo}. Additionally, a median ensemble filter of 21$\times$21 is applied on the obtained $\phi_{topo}^w$ following the processing steps provided by \cite{Cloude2005}.

\begin{equation}
    \label{eq:topo}
    \phi_{topo}^w = \arg\left(\widetilde{\gamma}\left(\vv{w_v}\right) - \widetilde{\gamma}\left(\vv{w_s}\right)\left(1 - L_{\vv{w_s}}\right)\right)
\end{equation}
where,
\begin{align*}
    & L_{\vv{w_s}} = \frac{-B - \sqrt{B^2 - 4AC}}{2A}, L_{\vv{w_s}} \in [0, 1] \\
    & A = \left\lvert\widetilde{\gamma}\left(\vv{w_v}\right)\right\rvert^2 - 1 \\
    & B = 2\Re\left(\widetilde{\gamma}\left(\vv{w_v}\right) - \widetilde{\gamma}\left(\vv{w_s}\right)\widetilde{\gamma}^*\left(\vv{w_v}\right)\right) \\
    & C = \left\lvert\widetilde{\gamma}\left(\vv{w_v}\right) - \widetilde{\gamma}\left(\vv{w_s}\right)\right\rvert^2
\end{align*}

Eventually, the SSD ($\Delta{Z_s}$) and SSWE ($=\rho_s\Delta{Z_s}$) are estimated wherein the standing snow density ($\rho_s = $ 0.315 g/cm\textsuperscript{3}) measured by the Dhundi SPA at 06:22 hrs IST on January 8, 2016, has been used. Here, $\eta = $ 0.65, the volume coherence threshold, $\tau_v = $ 0.6 (pixels having $\tau_v < $ 0.6 are neglected $\forall\tau_v \in$ [0, 1]), and the SSD values are averaged based on a 57$\times$57 ensemble filter window. The entire Pol-InSAR workflow is summarised in Figure \ref{fig:ssd_method} which shows the main processing blocks.

\begin{figure}[htb]
    \centering
    \includegraphics[width=\textwidth]{Figures/Methods/SSD_Method.png}
    \caption{SSD and SSWE estimation workflow using Pol-InSAR.}
    \label{fig:ssd_method}
\end{figure}

However, in order to compute the inverse $\sinc_\pi$ (normalised $\sinc$) function in Eq. \eqref{seq:hybrid_dem}, the \cite{Cloude2010} approximation ($\sinc_C^{-1}$) in Eq. \eqref{seq:un_cloude} is replaced by Eq. \eqref{seq:secant} where the secant method \citep{Cheney2012} has been applied to find $\alpha_r \in \mathbb{R}$ (rad), the desired root or inverse. Moreover, to make the \cite{Cloude2010} approximation compliant with the scientific computing libraries such as SciPy \citep{Jones2001} which use the $\sinc_\pi$ function, the normalised variant of Eq. \eqref{seq:un_cloude} given by Eq. \eqref{seq:ncloude} is incorporated where $\sinc_{\pi_C}^{-1}$ denotes the inverse of the $\sinc_\pi$ function computed using the \cite{Cloude2010} approach. Similarly, $\sinc_{\pi_S}^{-1}$ represents the inverse of the $\sinc_\pi$ function obtained by applying the secant method \citep{Cheney2012, Jones2001}. This root finding technique has been deployed as it is more accurate than the given approximation in Eq. \eqref{seq:ncloude}, the analysis of which is described in section \ref{sec:res}. Still, in the Python implementation, this approximation is used as an initial guess to the secant method for faster convergence. It is also used as a fallback option if the secant method is unable to converge within 50 iterations or the default convergence threshold of 1.4E-8 \citep{Jones2001}.

\begin{subequations}
    \centering
    \begin{align}
        \label{seq:un_cloude}
        & \sinc_C^{-1}\left(\gamma\left(\vv{w_v}\right)\right) = \pi - 2\sin^{-1}\left(\gamma\left(\vv{w_v}\right)^{0.8}\right) \\
        \label{seq:secant}
        & \sinc_\pi\alpha_r  - \gamma\left(\vv{w_v}\right) = 0 \\
        \label{seq:ncloude}
        & \sinc_{\pi_C}^{-1}\left(\gamma\left(\vv{w_v}\right)\right) = \frac{\sinc_C^{-1}\left(\gamma\left(\vv{w_v}\right)\right)}{\pi}
    \end{align}
\end{subequations}

\subsection{Validation, Uncertainty Assessment, and Sensitivity Analysis}
\label{ssec:vus}
\subsubsection{Validation Process}
\label{sssec:val}

One of the significant challenges in this work has been the limited ground-truth data availability. Since, in-situ data from only two ground stations are available, the conventional way of accuracy assessment through regression plots \citep{Kugler2015, Leinss2014, Kumar2017} is infeasible in this context. Moreover, the Kothi AWS area falls in the layover region for the descending pass acquisitions and hence, only the Dhundi region which is free from layover, shadow and foreshortening effects, is used for validation. In this case, a neighbourhood window of size 3$\times$3 ($\approx$ 81 m\textsuperscript{2} ground area) surrounding the Dhundi SPA is selected for validating the processed SD and SWE results by considering only the statistical mean and standard deviation.

\subsubsection{Uncertainty Assessment}
\label{sssec:ua}

Due to the complex terrain characteristics there exist significant uncertainty sources which could potentially lead to the overall degradation of the output accuracy. Although the winter-time data (January 8, 2016) is fully polarimetric, the summer-time data (June 8, 2017) is dual-pol. So, in order to comparatively understand the backscattering mechanisms in these two scenes, the dual-pol entropy ($H \in$ [0, 1]) and the scattering alpha angle ($\alpha \in$ [0$^\circ$, 90$^\circ$]) or $H/{\alpha}$ decomposition (5$\times$5 window size) is used \citep{Cloude2010, Lee2009, Singh2014}. This is carried out through the $H/{\alpha}$ plane plot which demarcates eight feasible zones (Z9 being the unclassified pixels) based on the different scattering classes as shown in Figure \ref{fig:ha}. It should be noted that, this diagram which follows the SNAP style \citep{ESA2018}, uses slightly different labels as compared to the \cite{Lee2009} $H/{\alpha}$ plane convention where
the labels Z1, Z2, Z3 are denoted as Z7, Z8, Z9 and vice-versa respectively. However, the scattering mechanisms are exactly the same in both these conventions. Here, the blue curve acts as a boundary to the plane which essentially denotes the reliability of the classification in high entropy conditions \citep{Brunner2009}.

\begin{figure}[htb]
    \centering
    \includegraphics[width=\textwidth]{Figures/Methods/HA.png}
    \caption{$H/{\alpha}$ plane showing different scattering zones. Z1: Dihedral, Z2: Dipole, Z3: Bragg Surface, Z4: Double
bounce, Z5: Anisotropic, Z6: Random surface, Z7: Complex structures, Z8: Random anisotropic, Z9: Non-feasible.}
    \label{fig:ha}
\end{figure}

The dual-pol $H/{\alpha}$ decomposition is further used by the unsupervised Wishart classifier (ten iterations) which classifies the SAR data based on these scattering mechanisms and a quantitative estimate of the number of pixels in each such class can be obtained \citep{Cloude2010, Lee2009}. Therefore, by knowing the scattering properties, the terrain features present in the study area can be understood along with their changes during the snow season. In turn, these ground features which include rough surfaces, shrubs, boulders, and human settlements reduce the copolar coherence amplitude ($\gamma_c$) thereby leading to underestimated FSD results \citep{Leinss2014, Leinss2016}. In addition, the decrease in the Pol-InSAR surface coherence amplitude, $\left(\gamma\left(\vv{w_v}\right) = \left\lvert\widetilde{\gamma}\left(\vv{w_v}\right)\right\rvert\right)$ may result in overestimated volumetric height \citep{Cloude2010, Hajnsek2009, Kugler2015}, in this case, SSD. Thus, the uncertainty assessment by means of the identification of the backscattering mechanisms constitutes a key role in this work.

Apart from this, the forest cover map (obtained from WRD, IIRS) along with the layover and shadow regions computed using SAR simulation are used to mask out the pixels which degrade the quality of the results. This is a standard approach used in the studies focusing on snow property estimation in forested or alpine terrains \citep{Leinss2014, Leinss2016, Singh2017, Thakur2012, Usami2016}.

\subsubsection{Sensitivity Analysis}
\label{sssec:sa}

The variation of the SD and SWE values corresponding to the changes in the free parameters in the FSD and SSD inversion models (window size, coherence threshold, scaling factors) are observed by iteratively running the algorithms and computing the statistical mean and standard deviation using the neighbourhood window discussed earlier in section \ref{sssec:val}. This helps in deciding the window shape and sizes and also choosing the optimum values for the several free parameters. Moreover, the accuracy of the root finding algorithm discussed in section \ref{ssec:ssd} is also checked for some possible coherence values.

In addition, the ground elevation measurements acquired during the field visit to Dhundi and Kothi have been compared with the ALOS PALSAR DEM elevations ($z$). The effect of the DEM errors on the LIA, $\theta_l$, is then checked for performing the sensitivity analysis (SA) using Eq. \eqref{eq:sa} which incorporates the slope angles in $x$ ($\omega_x$) and $y$ ($\omega_y$) directions (pixel co-ordinate system where $z$ is the corresponding elevation value) derived from the DEM elevation values along with the radar incidence angle ($\theta$) \citep{Lee2000, Lee2009}. Here, the terms $\mathrm{d}z/\mathrm{d}x$ and $\mathrm{d}z/\mathrm{d}y$ refer to the rate of elevation ($z$) change in the $x$ and $y$ directions respectively.

\begin{equation}
    \centering
    \label{eq:sa}
    \theta_l = \cos^{-1}\frac{\cos\omega_x\cos\left(\omega_y - \theta\right)}{\sqrt{\cos^2\omega_y\sin^2\omega_x + \cos^2\omega_x}}
\end{equation}
where, 
\begin{equation*}
    \centering
    \omega_x = \tan^{-1}\frac{\mathrm{d}z}{\mathrm{d}x}, \omega_y = \tan^{-1}\frac{\mathrm{d}z}{\mathrm{d}y}
\end{equation*}

\section{Results and Discussion}
\label{sec:res}
\subsection{Scattering Mechanisms}
\label{ssec:scat}

The winter (January 8, 2016) and summer-time (June 8, 2017) dual-pol $H/{\alpha}$ decomposition (Figure \ref{fig:ha_res}) and unsupervised Wishart classification (Figure \ref{fig:wishart_res}) results combined with the derived class percentage statistics (Figure \ref{fig:percent}) show that, in the presence of snow, the high entropy anisotropic volume scattering (Z8) increases by 5.11\% whereas the medium entropy volume scattering (Z5) decreases by 7.01\% for the entire study area. This reduction in the Z5 volume scattering could be attributed to the partially snow covered forests and shrubs which exhibit higher volume scattering at X-band during the snow-free season (Figure \ref{subfig:mountain}). The corresponding dual-pol Wishart classified maps are displayed along with the zoomed views in Figure \ref{subfig:wishart_jan} and Figure \ref{subfig:wishart_jun} respectively.

\begin{figure}[htb]
    \centering
    \begin{subfigure}[t]{\textwidth}
        \raisebox{-\height}{\includegraphics[width=\textwidth]{Figures/Results/wishart_jan.png}}
        \caption{}
        \label{subfig:wishart_jan}
    \end{subfigure}
    %%%%%%%%%%%%%%%%%%%%%%%%%%%%%%%%%%%second row
    \begin{subfigure}[t]{\textwidth}
        \raisebox{-\height}{\includegraphics[width=\textwidth]{Figures/Results/wishart_jun.png}}
        \caption{}
        \label{subfig:wishart_jun}
    \end{subfigure}
    \caption{Zoomed views over Dhundi of the Wishart classified maps for the \subref{subfig:wishart_jan} January 8, 2016, and \subref{subfig:wishart_jun} June 8, 2017 data. Here, only the layover and shadow mask has been applied. Also, the Kothi area is excluded from the analysis since it lies in the layover region.}
    \label{fig:wishart_res}
\end{figure}

Moreover, the Bragg surface scattering (Z3) is slightly higher in summer (10.88\%) as compared to the winter (10.38\%). One plausible reason for this is the 20 mm rainfall which occurred on June 7, 2017, evening (data retrieved from the Dhundi record book). Also, the occurrence of fresh snowfall in areas which did not have prior standing or old snow could result in surface scattering from the ground \citep{Leinss2014}. Apart from this, the asbestos gable roofs used in the human settlements (Figure \ref{subfig:base} and Figure \ref{subfig:landscape}) are strong single-bounce surface scatterers \citep{Brunner2009}. However, with snow accumulation on these materials, the surface scattering could be reduced. Another prominent feature noticeable in Figure \ref{subfig:wishart_jun} is the high amount of surface scattering from the river bed (Figure \ref{subfig:beas}) during the summer season. This is caused by both the boulders and the increasing flow of snow-melt water in the river (Figure \ref{subfig:beas}).

Furthermore, the human settlements result in double-bounce scattering (Z4) \citep{Brunner2009}, which in the winter-time scenario reduces by 0.34\%. Also, the random surface scattering (Z6) increases by 0.66\% which could be caused by the presence of small snow patches on the ground. Other than this, there is a strong decrease in the low entropy multiple (dihedral) scattering from 8.23\% to 5.17\% in the snow-covered season which could be caused by the added snow layer on the buildings and also boulders.

\begin{figure}[htb]
    \centering
    \includegraphics[width=\textwidth]{Figures/Results/Percent.png}
    \caption{Scattering class percentages (rounded to 2 decimal places) from the unsupervised Wishart classification. The different zone labels are described in Figure \ref{fig:ha}.}
    \label{fig:percent}
\end{figure}

Another interesting aspect in this context is the increase (from 9.93\% to 19.8\%) in the number of unclassified or non-feasible pixels (Z9) for the winter-time image (Figure \ref{fig:percent}) which is also depicted through the $H/{\alpha}$ plane plots in Figure \ref{subfig:ha_jan} and Figure \ref{subfig:ha_jun}. This is primarily resulting from the added terrain complexity owing to the snow accumulation. In order to resolve this issue, the quad-pol entropy ($H$), anisotropy ($A \in$ [0, 1]), alpha ($\alpha$), $H/A/{\alpha}$ decomposition has been applied on the January 8, 2016 data. The corresponding $H/{\alpha}$ plane plot in Figure \ref{subfig:ha_jan_quad} shows that the quad-pol approach is able to fully classify the winter-time image. However, since the summer-time image is having only HH and VV channels, the dual-pol method has been used to properly compare the respective scattering mechanisms \citep{Majumdar2019}.

\begin{figure}[htb]
    \centering
    \begin{subfigure}[t]{0.49\textwidth}
        \raisebox{-\height}{\includegraphics[width=\textwidth]{Figures/Results/ha_jan.png}}
        \caption{}
        \label{subfig:ha_jan}
    \end{subfigure}
    \hfill
    \begin{subfigure}[t]{0.49\textwidth}
        \raisebox{-\height}{\includegraphics[width=\textwidth]{Figures/Results/ha_jun.png}}
        \caption{}
        \label{subfig:ha_jun}
    \end{subfigure}
    %%%%%%%%%%%%%%%%%%%%%%%%%%%%%%%%%%%second row
    \begin{subfigure}[t]{0.49\textwidth}
        \raisebox{-\height}{\includegraphics[width=\textwidth]{Figures/Results/ha_jan_quad.png}}
        \caption{}
        \label{subfig:ha_jan_quad}
    \end{subfigure}
    \caption{Dual-pol $H/{\alpha}$ plane plots for the \subref{subfig:ha_jan} January 8, 2016, and \subref{subfig:ha_jun} June 8, 2017 data, \subref{subfig:ha_jan_quad} Quad-pol $H/{\alpha}$ plane plot for the January 8, 2016 data. The colours red, green, and blue indicate the point density with red being the highest, and blue as the lowest. These plots have been made using SNAP \citep{ESA2018}.}
    \label{fig:ha_res}
\end{figure}

Thus, from this discussion, it is clearly observed that the presence of snow causes a substantial change of the scattering patterns in the study area resulting in significant uncertainty sources. In turn, the optimisation of the model parameters along with the sensitivity analysis of the FSD and SSD values depend on these scattering types. As an example, if there is low surface scattering then the FSD inversion model leads to underestimated values \citep{Leinss2014} whereas for low volume scattering, the SSD results are generally underestimated \citep{Cloude2005, Hajnsek2009, Kugler2015}. Therefore, the uncertainty assessment by means of the scattering mechanism classification is one of the key aspects of this research.

\subsection{Sensitivity Analysis Results}
\label{ssec:sar}
In order to perform the sensitivity analysis, only the Dhundi area is chosen as Kothi is lying in the layover zone for this acquisition (January 8, 2016).

\subsubsection{FSD Model Parameters}
\label{sssec:fsd_param}
The FSD inversion model discussed in section \ref{ssec:fsd} apply the ensemble averaging operation twice--- once on the computed CPD and then subsequently on the output FSD values. As a result, the selection of an optimal window size in both these cases is critical in obtaining reliable estimates. At first, the mean ($\mu_{\gamma_c}$) and standard deviation ($\sigma_{\gamma_c}$) of the copolar coherence amplitude ($\gamma_c$) are checked over the Dhundi area based on a 3$\times$3 neighbourhood window (same as the validation window in section \ref{sssec:val}). The ensemble window for which the maximum $\mu_{\gamma_c}$ occurs is subsequently used for estimating the FSD following the methodology described in section \ref{ssec:fsd}. This selection procedure concerning the maximisation of $\mu_{\gamma_c}$ is depicted in Figure \ref{fig:fsd_coh} wherein the ensemble window of size 3$\times$3 is found to be suitable even though $\sigma_{\gamma_c} \approx$ 0.06 of this window is slightly higher than that ($\sigma_{\gamma_c} \approx$ 0.02) of the 5$\times$5 window.

\begin{figure}[htb]
    \centering
    \includegraphics[width=\textwidth]{Figures/Results/Coh_FSD.png}
    \caption{Effect of the window size on the mean and standard deviation of the copolar coherence amplitude. All the
values are rounded to 2 decimal places. Here, only odd window sizes are considered because these have been previously used in prior studies \citep{Leinss2018, Leinss2014, Kumar2011}.}
    \label{fig:fsd_coh}
\end{figure}

Next, the sensitivity of the FSD values with respect to the ensemble window size is taken into account. This is shown in Figure \ref{fig:fsd_res} where the analysis starts from the window size 3$\times$3 and continues till 99$\times$99 with an increment of two pixels in each direction (Figures \ref{subfig:fwindow_mean} and \ref{subfig:fwindow_var}). In this context, smaller window sizes ($<$ 45$\times$45) are not considered following the work of \cite{Leinss2014} where a 45$\times$35 window size has been chosen. Similarly, higher window sizes ($>$ 65$\times$65) are not used because of the varying topography in the study area. Another reason is that, since the Dhundi region exhibits moderate undulating terrains, so from the validation perspective, window sizes which cover ground areas of more than 0.4 km\textsuperscript{2} are excluded from the analysis. Hence, the transect as shown in Figure \ref{subfig:fwindow_transect} is used for deciding the FSD ensemble window size.

In this regard, the SPA measured FSD ground-truth data at 06:22 hrs January 8, 2016 (IST) was 18.7 cm, and that of the manual record book was 18 cm (section \ref{ssec:data}). Assuming the SPA sensor bias to be 5 cm for the SD, the FSD ground observation of 18 cm is taken as the true value. Accordingly, it is observed from Figure \ref{subfig:fwindow_transect} that the 65$\times$65 window leads to the most accurate (94.83\%) mean FSD ($\mu_f \approx$ 18.93 cm) with a low FSD standard deviation ($\sigma_f \approx$ 0.1 cm). The corresponding mean FSWE ($\mu_{fs} \approx$ 13.25 mm) and FSWE standard deviation ($\sigma_{fs} \approx$ 0.07 mm) are also in concordance (94.84\% accuracy) with the ground-truth FSWE of 12.6 mm. However, it should be noted that the axial ratio for the fresh snow particle, $a_x/a_z$ = 1.5 is kept as an invariant throughout the entire FSD workflow (Figure \ref{fig:fsd_method}) and its SA has not been carried out.

Here, the FSD and $\gamma_c$ values are significantly influenced by the mixed scattering mechanisms exhibited by
the ground features (Figure \ref{subfig:wishart_jan}) which are being considered for the averaging operation. Moreover, the
underlying assumption of a smooth surface in the FSD inversion model does not hold for such rough terrains and consequently, $\gamma_c$ is reduced \citep{Leinss2014}. Therefore, the FSD SA concludes that even though a sufficiently reliable FSD estimate has been achieved in the Dhundi area, the window sizes need to be adequately adjusted for different multi-temporal SAR images acquired over the same region thereby leading to a more robust parameter optimisation process.

\begin{figure}[htb]
    \centering
    \begin{subfigure}[t]{0.49\textwidth}
        \raisebox{-\height}{\includegraphics[width=\textwidth]{Figures/Results/fsd_mean_window.png}}
        \caption{}
        \label{subfig:fwindow_mean}
    \end{subfigure}
    \hfill
    \begin{subfigure}[t]{0.49\textwidth}
        \raisebox{-\height}{\includegraphics[width=\textwidth]{Figures/Results/fsd_var_window.png}}
        \caption{}
        \label{subfig:fwindow_var}
    \end{subfigure}
    %%%%%%%%%%%%%%%%%%%%%%%%%%%%%%%%%%%second row
    \begin{subfigure}[t]{0.49\textwidth}
        \raisebox{-\height}{\includegraphics[width=\textwidth]{Figures/Results/FSD_Res.png}}
        \caption{}
        \label{subfig:fwindow_transect}
    \end{subfigure}
    \caption{Effect of the window size on the \subref{subfig:fwindow_mean} mean and \subref{subfig:fwindow_var} standard deviation of the FSD estimates. Here, the transect between the two red points (45$\times$45 and 65$\times$65) are collectively shown in \subref{subfig:fwindow_transect} wherein all the values are rounded to 2 decimal places.}
    \label{fig:fsd_res}
\end{figure}

\subsubsection{SSD Model Parameters}
\label{sssec:ssd_param}

The SSD inversion model as described from the implementation or methodological perspective in section \ref{ssec:ssd} incorporates several user-defined free parameters. Thus, it is necessary to conduct an appropriate SA for the hybrid Pol-InSAR based volumetric height (SSD) retrieval algorithm. Accordingly, the various model parameters and their optimisation are discussed below.

\subsubsection*{Volume and Surface Coherence Ensemble Window}

The ensemble windows corresponding to the number of looks (L) in Eq. (12) must be suitably chosen so as to maximise both the volume coherence amplitude, $\gamma\left(\vv{w_v}\right)$, and the surface coherence amplitude, $\gamma\left(\vv{w_s}\right)$. As a result, the SA for these window sizes is an important aspect of this work.

The effects of $L$ on the mean volume coherence amplitude, $\mu_{\gamma\left(\vv{w_v}\right)}$ and the mean surface coherence
amplitude, $\mu_{\gamma\left(\vv{w_s}\right)}$ which are measured by applying the same 3$\times$3 neighbourhood window over Dhundi (section \ref{ssec:ssd}) along with the respective standard deviations, $\sigma_{\gamma\left(\vv{w_v}\right)}$ and $\sigma_{\gamma\left(\vv{w_v}\right)}$, are displayed in Figure \ref{fig:ssd_coh}. 

\begin{figure}[htb]
    \centering
    \includegraphics[width=\textwidth]{Figures/Results/Coh_SSD.png}
    \caption{Effect of the number of looks ($L$) on the volume and surface coherence. All the values are rounded to 2
decimal places.}
    \label{fig:ssd_coh}
\end{figure}

It can be seen that for the executed test cases, with increasing $L$, there is a general decreasing trend for both these coherences. So, for the SSD estimation, $L =$ 3 is chosen even though \cite{Cloude2005} suggests the usage of higher values of $L$. This is because, $\sigma_{\gamma\left(\vv{w_v}\right)} \approx$ 0.1 and $\sigma_{\gamma\left(\vv{w_s}\right)} \approx$ 0.18 are sufficiently small with adequately high $\mu_{\gamma\left(\vv{w_v}\right)} \approx$ 0.67 and $\mu_{\gamma\left(\vv{w_s}\right)} \approx$ 0.68. Also, since there is only one validation point for the entire study area, $L =$ 3 is justifiable.

However, there exist several free parameters in this Pol-InSAR based SSD inversion model (section \ref{ssec:ssd}) and hence, the volume and surface coherence ensemble windows need to be kept constant ($L =$ 3) for the subsequent SA of the other parameters.

\subsubsection*{Scaling Parameters}

It has been previously discussed in section \ref{ssec:ssd} that there are two scaling parameters involved in the SSD estimation process. These are the vertical wavenumber scaling parameter ($\eta^\prime\in\mathbb{R}_{>0}^+$) and the scaling factor ($\eta\in$ [0, 1]) of the hybrid DEM differencing approach developed by \cite{Cloude2010}. Here, the SA of only $\eta$ is carried out and $\eta^\prime = $ 5 is kept constant throughout the entire workflow. Also, the volume coherence threshold, $\tau_v = $ 0.6, $L =$ 3, ground phase median ensemble filter window (21$\times$21), vertical wavenumber ensemble average window (21$\times$21), and the SSD ensemble average window of size 57$\times$57 are unchanged during this SA.

The monotonically increasing SSD with respect to increasing $\eta$ are displayed in Figure \ref{fig:eta}. For $\eta =$ 0, the standard DEM differencing technique \citep{Cloude2005} results in the mean SSD, $\mu_s \approx$ 42.46 cm with the corresponding SSD standard deviation, $\sigma_s \approx$ 0.49 cm. As the SPA measured SSD at 06:22 hrs IST, January 8, 2016, is 54.9 cm, so $\mu_s$ is underestimated. Naturally, the mean SSWE, $\mu_{ss} \approx$ 133.76 mm (with SSWE standard deviation, $\sigma_{ss} \approx$ 1.53 mm) is also lower compared to the SPA measured SSWE of 173 mm. Thus, to effectively optimise the volumetric height, $\eta$ needs to be suitably increased \citep{Cloude2005, Cloude2010}.

\begin{figure}[htb]
    \centering
    \includegraphics[width=\textwidth]{Figures/Results/Eta.png}
    \caption{Increasing mean SSD with respect to the scaling parameter $\eta$.}
    \label{fig:eta}
\end{figure}

In this context, \cite{Cloude2005} has suggested setting $\eta = $ 0.4 for which the accuracy of the estimated tree height is found to be more than 90\%. Although by keeping $\eta = $ 0.4, $\mu_s \approx$ 49.64 cm ($\sigma_s \approx$ 0.54 cm) is obtained with $\sim$90.42\% accuracy, the complexity of the snow microstructure, anisotropy, and length scales necessitates the need for achieving even higher accuracies \citep{Leinss2016}. Moreover, in the presence of significantly varying hydrometeorological conditions which include high surface roughness and associated uncertainty sources (section \ref{ssec:scat}), the volume and surface coherence amplitudes generally do not reach expected values of higher than 0.8 \citep{Cloude2005, Kugler2015}. Therefore, with $\eta = $ 0.65, the best SSD and SSWE accuracies of 99.53\% ($\mu_s \approx$ 54.64 cm) and 99.48\% ($\mu_{ss} \approx$ 172.10 mm) respectively are achieved over Dhundi with low standard deviations ($\sigma_s \approx$ 0.58 cm, $\sigma_{ss} \approx$ 1.82 mm) accounting for high reliability. These results highlight the significance of this scaling parameter $\eta$ towards controlling the snow structural height variations \citep{Cloude2005, Cloude2010} and hence, the robustness of the hybrid DEM differencing model (section \ref{ssec:ssd}) is verified.

\subsubsection*{Computing $\sinc$ Inverse}
In order to test the accuracy of the $\sinc_\pi$ inverse function, sample test data representing the actual inverse, $\alpha_r$, have been prepared as shown in Table \ref{table:2}. Next, the $\sinc_\pi$ of these data, $\sinc_{\pi}\left(\alpha_r\right)$, is computed which essentially corresponds to the possible $\gamma\left(\vv{w_v}\right)$ values. So, the idea of performing SA in this scenario is to check the accuracy of the calculated $\sinc_{\pi_C}^{-1}$ (Eq. \eqref{seq:ncloude}) and $\sinc_{\pi_S}^{-1}$ (Eq. \eqref{seq:secant}) of the $\sinc_{\pi}\left(\alpha_r\right)$ values by comparing these with $\alpha_r$.

\begin{table}[ht]
\centering
\caption{Comparison between the normalised \cite{Cloude2010} sinc inverse and the secant sinc inverse methods.}
\label{table:2}
\begin{tabular}{c c c c}
\hline
\boldmath$\alpha_r$ \textbf{(rad)} & \boldmath$\sinc_{\pi}\left(\alpha_r\right)$   & \boldmath$\sinc_{\pi_C}^{-1}$ \textbf{(rad)}     & \boldmath$\sinc_{\pi_S}^{-1}$ \textbf{(rad)} \\ \hline
0.1                 & 0.984            & 0.103          & 0.100  \\ 
0.2                 & 0.935            & 0.206          & 0.200  \\ 
0.3                 & 0.858            & 0.308          & 0.300  \\ 
0.4                 & 0.757            & 0.409          & 0.400  \\ 
0.5                 & 0.637            & 0.509          & 0.500   \\ 
0.6                 & 0.505            & 0.607          & 0.600  \\ 
0.7                 & 0.368            & 0.703          & 0.700   \\ 
0.8                 & 0.234            & 0.798          & 0.800 \\ 
0.9                 & 0.109            & 0.891          & 0.900 \\ \hline

\end{tabular}
\end{table}

From Table \ref{table:2} it is observed that the secant method converges exactly (up to 13 decimal places) to the
actual $\alpha_r$ while the normalised \cite{Cloude2010} approximation of the $\sinc_\pi$ inverse has some minute errors
involved (RMSE $\approx$ 0.02 rad). Similarly, the $\sinc$ function is tested (Table \ref{table:3}) where $\sinc_C^{-1}$ and $\sinc_S^{-1}$ denote the standard \cite{Cloude2010} approximation (Eq. \eqref{seq:un_cloude}) and the secant method of root finding for the traditional $\sinc$ function respectively. Again, the secant method exactly converges (up to 13 decimal
places) whereas RMSE $\approx$ 0.02 rad is associated with the $\sinc_C^{-1}$. Here, the computed results shown in
Table \ref{table:2} and Table \ref{table:3} are rounded to 3 decimal places.

\begin{table}[ht]
\centering
\caption{Comparison between the traditional \cite{Cloude2010} sinc inverse and the secant sinc inverse methods.}
\label{table:3}
\begin{tabular}{c c c c}
\hline
\boldmath$\alpha_r$ \textbf{(rad)} & \boldmath$\sinc\left(\alpha_r\right)$   & \boldmath$\sinc_C^{-1}$ \textbf{(rad)}     & \boldmath$\sinc_S^{-1}$ \textbf{(rad)} \\ \hline
0.1                 & 0.998            & 0.103          & 0.100  \\ 
0.2                 & 0.993            & 0.207          & 0.200  \\ 
0.3                 & 0.985            & 0.310          & 0.300  \\ 
0.4                 & 0.974            & 0.413          & 0.400  \\ 
0.5                 & 0.959            & 0.516          & 0.500   \\ 
0.6                 & 0.941            & 0.618          & 0.600  \\ 
0.7                 & 0.920            & 0.721          & 0.700   \\ 
0.8                 & 0.897            & 0.823          & 0.800 \\ 
0.9                 & 0.870            & 0.925          & 0.900 \\ \hline

\end{tabular}
\end{table}

Therefore, by performing the SA of the $\sinc_{\pi_C}^{-1}$, $\sinc_{\pi_S}^{-1}$, $\sinc_C^{-1}$, and $\sinc_S^{-1}$, it is clearly understood that the secant method provides highly accurate results. Hence, in this work, $\sinc_{\pi_S}^{-1}$ is applied for solving Eq. \eqref{seq:hybrid_dem} wherein the $\sinc_{\pi_C}^{-1}\left(\gamma\left(\vv{w_v}\right)\right)$ value is used as an initial guess to the secant method for faster convergence.

\subsubsection*{SSD Ensemble Window}

Another essential free parameter used in the Pol-InSAR based SSD estimation model (section \ref{ssec:ssd}) is the SSD ensemble averaging window size. By keeping $\eta = $ 0.65, $\eta^\prime = $ 5, and other ensemble window sizes constant, the SA has been carried out to observe the SSD variations which are shown in Figure \ref{fig:swindow}. Here, the ensemble windows are the same which have been previously applied for the FSD values (section \ref{sssec:fsd_param}) so as to appropriately compare the SSD and FSD estimates.

\begin{figure}[htb]
    \centering
    \includegraphics[width=\textwidth]{Figures/Results/SSD_Res.png}
    \caption{Effect of the ensemble window size on the SSD values.}
    \label{fig:swindow}
\end{figure}

The graphical representation in Figure \ref{fig:swindow} shows that when the window size is increased beyond 57$\times$57, the SSD values increase sharply whereas, between the windows 53$\times$53 and 57$\times$57, the values are mostly similar. This could be attributed by the fact that, in mountainous terrains, elevation, and not distance, plays a critical role in controlling the snow accumulation \citep{Liu2017, Singh2014, Singh2017, Thakur2012}. The varying topographical conditions prominently visible in Figure \ref{fig:field} also ascertain that for larger window sizes, the snow depth variability could increase if a nearby mountain also lies within the neighbourhood window. So, considering these aspects, the ensemble window size of 57$\times$57 is selected which results in $\mu_s \approx$ 54.64 cm with $\sigma_s \approx$ 0.58 cm as discussed in the scaling parameter SA.

\subsubsection{DEM and LIA Error Analysis}
\label{sssec:error}

During the field visit (section \ref{sssec:field}), several DGPS points which had been acquired are used to check the
accuracy of the ALOS PALSAR DEM. In essence, the observed errors are then used to analyse the change in the LIA (Eq. \eqref{eq:sa}) induced by the corrected DEM (the erroneous DEM pixels are replaced by the respective DGPS measurements). 

The DEM errors calculated using the Dhundi and Kothi DGPS readings are displayed in Figure \ref{subfig:dem_error} and the subsequent LIA differences (computed from the corrected and original DEMs) for these points are shown in Figure \ref{subfig:lia_error}. As seen from these graphs, the absolute elevation errors range from 0.08 m to 16.30 m in the Dhundi region, whereas these vary from 0.19 m to 25.32 m in the Kothi area. Accordingly, the RMSE values for the elevation errors are approximately 6.71 m and 8.8 m respectively.

\begin{figure}[!htbp]
    \centering
    \begin{subfigure}[t]{\textwidth}
        \raisebox{-\height}{\includegraphics[width=\textwidth]{Figures/Results/DEM_Error.png}}
        \caption{}
        \label{subfig:dem_error}
    \end{subfigure}
    %%%%%%%%%%%%%%%%%%%%%%%%%%%%%%%%%%%second row
    \begin{subfigure}[t]{\textwidth}
        \raisebox{-\height}{\includegraphics[width=\textwidth]{Figures/Results/LIA_Error.png}}
        \caption{}
        \label{subfig:lia_error}
    \end{subfigure}
    \caption{\subref{subfig:dem_error} Absolute DEM errors obtained by comparing ALOS PALSAR DEM and the DGPS measurements and \subref{subfig:lia_error} observed absolute LIA errors. Here DB is the Dhundi base station point, D1-D86 are acquired in the Dhundi region, and K1-K72 are measured in the Kothi area using the DGPS.}
    \label{fig:error}
\end{figure}

In addition, the LIA varies from 0$^\circ$ to 7.59$^\circ$ (Dhundi) and 0$^\circ$ to 0.17$^\circ$ (Kothi) in these areas with the corresponding RMSE being nearly 2.54$^\circ$ and 0.02$^\circ$. However, since the LIA is dependent on the slope values (Eq. \eqref{eq:sa}), the DEM errors do not significantly influence the LIA. Also, in the FSD inversion model and the vertical wavenumber calculation (used in the SSD estimation) given by Eq. \eqref{eq:7} and Eq. \eqref{seq:vn} respectively, the sine ($\sin$) of the LIA is considered. So, the minute changes in the LIA do not strongly affect the FSD and SSD estimates which are obtained after applying sufficient ensemble averaging operation (section \ref{sec:method}). Evidently, the LIA only changes by about 1.9$^\circ$ near the Dhundi base station and thus, the FSD and SSD results are not exhibiting any sizeable impact from the associated DEM errors.

Therefore, the SA concerning the DEM errors and its propagation highlights that the subsequent LIA errors are not directly governed by the changes in the elevation values, rather the slopes in $x$ and $y$ directions (section \ref{sssec:sa}) act as the primary error sources. Also, the ALOS PALSAR DEM is sufficiently accurate even in the complex terrains and hence, its usage in the LIA computation is justified.

\subsection{Comparative Analysis of the Estimates}
\label{ssec:snow}

In order to visually observe the spatial patterns of the FSD and SSD estimates, suitable maps have been prepared which are shown in Figure \ref{subfig:fsd_map} and Figure \ref{subfig:ssd_map} respectively. Here, the resultant FSWE and SSWE maps are not shown as these have been computed by multiplying the constant snow densities, $\rho_f = $ 0.07 g/cm\textsuperscript{3} and $\rho_s = $ 0.315 g/cm\textsuperscript{3} to the FSD and SSD values respectively. Therefore, they have similar spatial characteristics like those of the snow depth maps.

\afterpage{\FloatBarrier}
\begin{figure}[htb]
    \centering
    \begin{subfigure}[t]{\textwidth}
        \raisebox{-\height}{\includegraphics[width=\textwidth]{Figures/Results/FSD.png}}
        \caption{}
        \label{subfig:fsd_map}
    \end{subfigure}
    %%%%%%%%%%%%%%%%%%%%%%%%%%%%%%%%%%%second row
    \begin{subfigure}[t]{\textwidth}
        \raisebox{-\height}{\includegraphics[width=\textwidth]{Figures/Results/SSD.png}}
        \caption{}
        \label{subfig:ssd_map}
    \end{subfigure}
    \caption{Zoomed views of the \subref{subfig:fsd_map} FSD map and \subref{subfig:ssd_map} SSD map for January 8, 2016. Here, the ground points surveyed (section \ref{sssec:field}) are shown wherein the closely spaced points have been acquired using the DGPS kinematic mode and fall on the nearby roads in the Dhundi region. The other points including the Dhundi base are measured using the static mode. Since the Kothi area falls in the layover and shadow zone, it is excluded from the zoomed view analysis.}
    \label{fig:sd_maps}
\end{figure}

As discussed earlier in sections \ref{sssec:fsd_param} and \ref{sssec:ssd_param}, the optimal FSD and SSD ensemble window sizes are 65$\times$65 and 57$\times$57 respectively. The maps in Figure \ref{fig:sd_maps} show these ensemble averaged estimates wherein $\mu_f \approx$ 18.93 $\pm$ 5.03 cm ($\sigma_f \approx$ 0.1 cm) and $\mu_s \approx$ 54.64 $\pm$ 5.19 cm ($\sigma_s \approx$ 0.58 cm) are observed over the 3$\times$3 neighbourhood window surrounding the Dhundi area with the corresponding $\mu_{fs} \approx$ 13.25 $\pm$ 5.02 mm ($\sigma_{fs} \approx$ 0.07 mm) and $\mu_{ss} \approx$ 172.10 $\pm$ 5.61 mm ($\sigma_{ss} \approx$ 1.82 mm). Here, the uncertainties are calculated based on the standard error of the estimate and the SPA measurement biases of 5 cm and 5 mm for the SD and SWE have been assumed respectively.

In addition, the overall mean FSD and SSD are 17.90 cm and 112.17 cm respectively wherein the standard deviations are found to be $\sim$6.46 cm and $\sim$30.80 cm. Accordingly, the mean FSWE and SSWE are $\sim$12.12 mm and $\sim$377.81 mm respectively where the associated standard deviations are $\sim$4.46 mm and $\sim$101.55 mm. 

The high SSD and SSWE standard deviations for the complete region highlight the extreme topographical conditions present in the study area. These variations can be confirmed from the ground survey (section \ref{sssec:field}) where the points (shown in Figure \ref{fig:sd_maps} had been acquired by considering the terrain undulations. Also, the aspect, slope, and elevation significantly influence the FSD and SSD estimates, the details of which have been discussed in the previous section.

Apart from this, it can be observed that these estimates are lower in the Dhundi base station area as compared to the surrounding regions. This phenomenon can be attributed to the presence of the human settlements (Figure \ref{subfig:base}) near the base point and are expected to have less snow accumulation than the natural surroundings. Moreover, the effect of multiple or double bounce scattering (Z4) near the Dhundi base is prominent even during the winter (Figure \ref{subfig:wishart_jan}). So, this could effectively reduce the copolar, volume and surface coherences (sections \ref{sssec:fsd_param} and \ref{sssec:ssd_param}) thereby explaining this observation.

\section{Conclusion and Future Scope}
\label{sec:conc}

The primary focus of this research lies in estimating the snow depth using which the snow water equivalent has been measured. Here, two different types of snow have been considered— freshly fallen (new) snow and standing (old) snow. In order to compute the FSD, the CPD method has been applied (section \ref{ssec:fsd}) on the January 8, 2016, TSX-TDX CoSSC bistatic dataset acquired over the Beas watershed, northwestern Himalayas (section \ref{sec:study}). Additionally, the hybrid DEM differencing and coherence amplitude inversion algorithm based on the single-baseline Pol-InSAR technique has been utilised to estimate the SSD for the same dataset (section \ref{ssec:ssd}). Also, the corresponding FSWE and SSWE are
obtained by multiplying constant fresh and standing snow densities.

Due to the complex hydrometeorological and topographical conditions of the study area (section \ref{sssec:geo}), significant uncertainty sources are present. These include the forests, boulders, highly rough surfaces, and human settlements (Figure \ref{fig:field}) which substantially reduce the surface and volume scattering coherences required to estimate the snow depths with adequate accuracy (\ref{ssec:vus}). Moreover, the limited ground-truth data availability has always been a major challenge from the onset of this work (section \ref{ssec:data}). Apart from this, the SAR data are affected by layover, shadowing and foreshortening in mountainous terrains and hence, these errors are inherently propagated through the subsequent processing steps. Furthermore, the PolSAR and Pol-InSAR models involve several user-defined parameters which have to be optimised (section \ref{sec:method}. In short, these are the main concerns involved in this work which are addressed by means of identifying the potential uncertainty sources ($H/{\alpha}$ decomposition and Wishart classification) and performing appropriate sensitivity analysis (section \ref{sssec:sa}).

Thus, the novelty of this research lies in suitably modifying and ultimately improving the hybrid Pol-InSAR model (section \ref{ssec:ssd}) to estimate the SSD which is new in the context of cryospheric studies. Additionally, the PolSAR CPD method for FSD retrieval has been tested for the first time in the presence of extreme topographically varying conditions. Although the FSD and SSD ground-truth measurements from only the Dhundi station had been available, the high accuracies of 94.83\% and 99.53\% respectively imply that these improved models work sufficiently well under the complex hydrometeorological situations.

As part of future work, it is recommended to use the multi-baseline Pol-InSAR technique \citep{Cloude2010} wherein $k_z$ can be simulated (instead of scaling by $\eta^\prime$) after an appropriate accuracy assessment \citep{Kumar2017}. Similarly, the effect of different window shapes (square or rectangular) and sizes can be considered for the ensemble averaging operation. This sort of sensitivity analysis will help in deciding optimal window structures separately for each model. Moreover, it is recommended to apply scattering mechanism based masks in conjunction with snow masks prepared from the high resolution optical datasets such as those provided by Sentinel-2 \citep{Zhu2015}. Additionally, the prior classification of the dry and wet snow including the preparation of snow cover maps \citep{Leinss2018, Thakur2012, Zhu2015} as necessary preprocessing steps will certainly improve the uncertainty assessment process.

Also, the use of the newer multi-temporal high resolution L-band datasets acquired by the upcoming SAR missions \citep{Krieger2016, Rosen2017} is recommended to further verify and validate these models. Moreover, radar altimeters such as the Ka-band InSAR altimeter could potentially improve the SD and SWE estimates, and could also be used for operational snow depth monitoring on a large-scale \citep{Hensley2016, Kim2018, Moller2011, Speziali2018}.

In this work, due to the time constraints, only one dataset (January 8, 2016) has been used for analysis. Instead, if a full scale time series analysis is performed, then the robustness of the SSD and FSD retrieval models can be checked. Furthermore, Pol-InSAR coherency optimisation can be carried out to suitably adjust the scattering phase centres\citep{Cloude2005, Cloude2010}. Moreover, the snow densities need to be computed gridwise (or if possible, pixelwise) by using hydrological modelling approaches \citep{Bartelt2002, Liang1994}. These can also be estimated from the PolSAR based techniques which are in practice \citep{Singh2017, Thakur2012}. Finally, appropriate statistical significance testing needs to be carried out to quantify further the uncertainties associated with the FSD and SSD estimates.

\section*{Acknowledgements}
This research work was carried out as part of the ISRO EOAM mountain ecosystem, TANDEM-X AO and ALOS-RA4 project (EOAM-ME (WRD)) on the Himalayan glaciers. Also, this work was conducted within the Indian Institute of Remote Sensing (IIRS), ISRO and University of Twente, Faculty ITC joint education programme (JEP) framework. Therefore, the authors are grateful to IIRS, ISRO and University of Twente, Faculty ITC along with the FOSS4G community for providing the necessary means to conduct this study.  
\section*{References}

\bibliography{refs}

\end{document}